\magnification \magstephalf
\nopagenumbers
\parindent = 0pt
\font\lgbf = cmbx10 scaled \magstep1
\centerline {\lgbf Course Information}
\smallskip
\centerline {\lgbf Mathematics 152, Calculus II, Honors Section}
\smallskip
\centerline {\lgbf Fall Semester 2004}
\bigskip
{\bf Homepage} 
\smallskip
http://www.rci.rutgers.edu/$\sim$yzhuang/math/152-f-04.html %\hfil\break
\bigskip
{\bf Weekly Schedule} 
\smallskip
\def \as {\phantom {$\ast$}}
\halign{#\hfil& \qquad\qquad #\hfil& \qquad\qquad \hfil#\hfil&
\qquad\qquad #\hfil\cr
Day& Component& Period& Location\cr
\noalign{\smallskip}
Monday& Lecture& 6& Hill 124 \cr
Wednesday& Lecture& 6& Hill 124\cr
Thursday& Workshop& 6& SEC 211 \cr}
\bigskip
{\bf Instructor} \enspace Yi-Zhi Huang %\hfil\break

\qquad {\bf Office} \enspace Hill 332 (Busch) %\hfil\break

\qquad {\bf Telephone} \enspace (732) 445-1314 %\hfil\break

\qquad {\bf Email} \enspace yzhuang@math.rutgers.edu %\hfil\break

\qquad {\bf Homepage} \enspace 
http://www.math.rutgers.edu/$\sim$yzhuang %\hfil\break

\qquad {\bf Office Hours} \enspace Monday and Wednesday, 2:30 to 3:30 pm
and by appointment
\bigskip
{\bf Peer Mentor} \enspace Jenner Yeh
\bigskip
{\bf Examination Dates}
\smallskip
\halign{#\hfil& \qquad\qquad#\hfil\cr
First Examination& Wednesday, October 13\cr
Second Examination& Monday, November 29\cr
Final Examination& Thursday, December 16, 4-7 PM\cr}
\smallskip
Both the first and second examinations are held at the regular class time and
in the regular meeeting room. The location of the final examination
will be announced several weeks before the end of the term.
\bigskip
{\bf Text} \enspace 
Calculus, Early Transcendentals, Stewart, $5^{\rm th}$ Edition,
Brooks-Cole Publishing Co. (This edition of the textbook differs in 
important ways from previous editions. Unfortunately the older editions 
will not be workable.) 
\bigskip
{\bf Graphing Calculator} \enspace 
A graphing calculator is required for this course. 
We have traditionally used the TI-82 or 83 and recommend either of them, 
but any calculator with equivalent capacities can be used, such as the 
popular TI-85 or 86. Calculators {\it will not be permitted on final exams}. 
\bigskip
{\bf Grading Policy} \enspace 
The various components of the course are weighted
as follows in the determination of your course grade.
\smallskip
\halign{\qquad #&  #\hfil& \qquad #\hfil\cr
&First hour exam:&              100 points\cr
&Second hour exam:&             100 points\cr
&Written workshop assignments:&  50 points\cr
&Quizzes and homework:&          50 points\cr
&Final exam:&                   200 points\cr
\noalign{\smallskip}
&Total:&                        500 points\cr}
\bigskip
{\bf Collection of Written Assignments and Attendance} %\hfil\break
\smallskip
1. Written workshop assignments are due the following Thursday in the
workshop. {\it Late write-ups will not be accepted.} \hfil\break
2. Similarly,  {\it late homework will
not be accepted}. \hfil\break		
3. Workshop attendance is mandatory and significant absence will
adversely affect your grade.
\bigskip
{\bf Examination Rules} %\hfil\break
\smallskip
No books, notes, or calculators may be used in taking
the examinations.

\bigskip
{\bf Schedule for Homework}
\smallskip
\halign{#\quad& #\hfil& \enspace #\hfil&\qquad #\hfil \cr
&Date& &Sections from which homework is due\cr
&Wednesday,   &September 8    &6.1, 6.2  \cr
&Monday,   &September 13   &6.3  \cr
&Monday,   &September 20  &6.4, 6.5, 7.1  \cr
&Monday,   &September 27      &7.2, 7.3, 7.4, 7.5  \cr
&Monday,   &October 4      &7.7, 7.8   \cr
&Monday,    &October 11     & 8.1, 8.2, 9.1, 9.2\cr
&Monday,   &October 25     &9.3, 9.4, 11.1\cr
&Monday,   &November 1    &11.2, 11.3, 11.4   \cr
&Monday,   &November 8     &11.5, 11.6, 11.7   \cr
&Monday,   &November 15   &11.8, 11.9   \cr
&Monday,   &November 22    &11.10, 11.11   \cr
&Monday,   &December 6     &11.12 \cr
&Monday,   &December 13    &10. 1, 10.2, 10.3, 10.5   \cr}
\bigskip
{\bf Workshop: Homework, Quizzes, Workshops} %\hfil\break
\smallskip
The Thursday class meeting will be devoted to going over homework problems,
having short quizzes, and doing workshop problems.
\smallskip
The first half hour of the Workshop period will be reserved 
for going over homework. Sometimes a short quiz may be given. The
remainder of the time will be used for doing workshop problems.
These problems will consist generally of 3-6 problems, handed out at
the beginning of the workshop.  The workshop problems are generally
somewhat more difficult and open-ended than the regular homework
problems. Students work on these problems in small groups and 
cooperative effort is encouraged. While joint work is appropriate
for the workshop, the final write-up of the problems should be your
own. The instructor and peer mentor are there to advise you 
with strategies for approaching the problem. In 
the workshops devoted to review exams for the hour exams and for the
final exam, they will answer any questions that you
want answered. %\hfil\break
\smallskip
The workshop problems that will be collected will be announced at
the end of each workshop.  These are due during the following
workshop.  The grading will take into account not only the accuracy
of your solution, but also the quality of your exposition. Thus, for
example, the problem should be stated, all notation defined, diagrams
clearly labelled, and steps fully explained. Neatness and legibility
are important.
\smallskip
Each week homework problems will be collected. 
Again, presentation, as well as mathematical
correctness, is important.
\bigskip
{\bf Final Examination} \hfil \break
The final examination given in this course is common to all sections
of Mathematics 152.

 \vfil \eject \bye







